% GENERATED FILE. Edit canonical lessons, then run npm run generate:books.
% canonical-source-hash: fnv1a64:e83b299117a3dbb3
% canonical-lessons: ES-C712-sensacion, ES-C712-animo, ES-C712-desanimado, ES-C712-animar, ES-C712-alivio

\chapter{Mi estado de ánimo}
\label{ch:es-mi-estado-de-animo}
\hlchaptermodality{712}

\begin{chapteropening}
\textbf{By the end of this chapter:} \emph{I can describe my mood and cheer up a friend who is down.}

Complete the last lesson of chapter 712: I can describe my mood and cheer up a friend who is down.
\end{chapteropening}

\section[la sensación]{la sensación — the sensation, the feeling}

\label{lesson:ES-C712-sensacion}



\begin{quote}
Before the new one: say the Spanish for moving, then the Spanish for to move.
\end{quote}

\subsection*{You'll want to know: la sensación}
\textbf{la sensación} — \textquotedblleft{}the sensation, the feeling\textquotedblright{}.

\begin{quote}
Tengo la sensación de que va a llover.
\end{quote}

\begin{quote}
\emph{I have a feeling it is going to rain.}
\end{quote}

\emph{Tengo la sensación de que...} — I have the feeling that...

\begin{cousinweb}
\emph{sensación} is usually described as a learned borrowing from Medieval Latin \emph{sēnsātiō}, built on \emph{sēnsus} \textquotedblleft{}feeling\textquotedblright{}, from \emph{sentīre} \textquotedblleft{}to feel\textquotedblright{}.
\end{cousinweb}

\begin{grammarlens}[title={What you've built}]
One more word for telling, explaining and describing at B1.
\end{grammarlens}

\subsection*{Your turn}
\begin{itemize}
\raggedright
  \item \emph{Say it:} \emph{la sensación}
  \item \emph{Say it:} \emph{la sensación}, once more
  \item \emph{Say it:} a sentence of your own with \emph{la sensación}
  \item \emph{Recall:} say the Spanish for moving, then the Spanish for to move, then say \emph{la sensación} again
\end{itemize}

\begin{tcolorbox}[breakable,colback=teal!4,colframe=teal!35!black,title={Before you move on}]
What does la sensación mean? (\textquotedblleft{}The sensation, the feeling\textquotedblright{}.) Say it once more.
\end{tcolorbox}
\section[el ánimo]{el ánimo — mood, spirits; courage}

\label{lesson:ES-C712-animo}



\begin{quote}
Before the new one: say the Spanish for to move, then the Spanish for the sensation.
\end{quote}

\subsection*{You'll want to know: el ánimo}
\textbf{el ánimo} — \textquotedblleft{}mood, spirits; courage\textquotedblright{}.

\begin{quote}
¡Ánimo! Tú puedes.
\end{quote}

\begin{quote}
\emph{Come on! You can do it.}
\end{quote}

\emph{¡Ánimo!} is what you shout to encourage someone. \emph{Estado de ánimo} — mood.

\begin{cousinweb}
\emph{ánimo} is a learned borrowing from Latin \emph{animus} \textquotedblleft{}spirit, mind\textquotedblright{} (EtymDB); English \emph{unanimous}, of one mind, is built on it.
\end{cousinweb}

\begin{grammarlens}[title={What you've built}]
One more word for telling, explaining and describing at B1.
\end{grammarlens}

\subsection*{Your turn}
\begin{itemize}
\raggedright
  \item \emph{Say it:} \emph{el ánimo}
  \item \emph{Say it:} \emph{el ánimo}, once more
  \item \emph{Say it:} a sentence of your own with \emph{el ánimo}
  \item \emph{Recall:} say the Spanish for to move, then the Spanish for the sensation, then say \emph{el ánimo} again
\end{itemize}

\begin{tcolorbox}[breakable,colback=teal!4,colframe=teal!35!black,title={Before you move on}]
What does el ánimo mean? (\textquotedblleft{}Mood, spirits; courage\textquotedblright{}.) Say it once more.
\end{tcolorbox}
\section[desanimado]{desanimado — discouraged, down}

\label{lesson:ES-C712-desanimado}



\begin{quote}
Before the new one: say the Spanish for the sensation, then the Spanish for mood.
\end{quote}

\subsection*{You'll want to know: desanimado}
\textbf{desanimado} — \textquotedblleft{}discouraged, down\textquotedblright{}.

\begin{quote}
Está desanimado porque no encuentra trabajo.
\end{quote}

\begin{quote}
\emph{He is down because he cannot find work.}
\end{quote}

Used with \emph{estar}. The opposite is \emph{animado}.

\begin{cousinweb}
\emph{desanimado} is built on \emph{animado} with the prefix \emph{des-}; \emph{animar} comes from Latin \emph{animāre} \textquotedblleft{}to give life\textquotedblright{}, from \emph{anima} \textquotedblleft{}breath, soul\textquotedblright{} (EtymDB).
\end{cousinweb}

\begin{grammarlens}[title={What you've built}]
One more word for telling, explaining and describing at B1.
\end{grammarlens}

\subsection*{Your turn}
\begin{itemize}
\raggedright
  \item \emph{Say it:} \emph{desanimado}
  \item \emph{Say it:} \emph{desanimado}, once more
  \item \emph{Say it:} a sentence of your own with \emph{desanimado}
  \item \emph{Recall:} say the Spanish for the sensation, then the Spanish for mood, then say \emph{desanimado} again
\end{itemize}

\begin{tcolorbox}[breakable,colback=teal!4,colframe=teal!35!black,title={Before you move on}]
What does desanimado mean? (\textquotedblleft{}Discouraged, down\textquotedblright{}.) Say it once more.
\end{tcolorbox}
\section[animar]{animar — to cheer up, to encourage}

\label{lesson:ES-C712-animar}



\begin{quote}
Before the new one: say the Spanish for mood, then the Spanish for discouraged.
\end{quote}

\subsection*{You'll want to know: animar}
\textbf{animar} — \textquotedblleft{}to cheer up, to encourage\textquotedblright{}.

\begin{quote}
Mis amigos me animan a seguir.
\end{quote}

\begin{quote}
\emph{My friends encourage me to keep going.}
\end{quote}

\emph{Animar a alguien a} + infinitive. \emph{Animarse} — to cheer up, to decide to do something.

\begin{cousinweb}
\emph{animar} is inherited from Latin \emph{animāre} \textquotedblleft{}to fill with breath, to give life\textquotedblright{} (EtymDB), from \emph{anima} \textquotedblleft{}breath, soul\textquotedblright{}.
\end{cousinweb}

\begin{grammarlens}[title={What you've built}]
One more word for telling, explaining and describing at B1.
\end{grammarlens}

\subsection*{Your turn}
\begin{itemize}
\raggedright
  \item \emph{Say it:} \emph{animar}
  \item \emph{Say it:} \emph{animar}, once more
  \item \emph{Say it:} a sentence of your own with \emph{animar}
  \item \emph{Recall:} say the Spanish for mood, then the Spanish for discouraged, then say \emph{animar} again
\end{itemize}

\begin{tcolorbox}[breakable,colback=teal!4,colframe=teal!35!black,title={Before you move on}]
What does animar mean? (\textquotedblleft{}To cheer up, to encourage\textquotedblright{}.) Say it once more.
\end{tcolorbox}
\section[el alivio]{el alivio — relief}

\label{lesson:ES-C712-alivio}



\begin{quote}
Before the new one: say the Spanish for discouraged, then the Spanish for to cheer up.
\end{quote}

\subsection*{You'll want to know: el alivio}
\textbf{el alivio} — \textquotedblleft{}relief\textquotedblright{}.

\begin{quote}
¡Qué alivio! Ya llegaste.
\end{quote}

\begin{quote}
\emph{What a relief! You are here.}
\end{quote}

\emph{¡Qué alivio!} — What a relief!

\begin{cousinweb}
\emph{alivio} is built on \emph{aliviar}, which comes from Latin \emph{allevāre} or \emph{alleviāre} \textquotedblleft{}to lighten\textquotedblright{} (EtymDB), from \emph{levis} \textquotedblleft{}light\textquotedblright{}.
\end{cousinweb}

\begin{grammarlens}[title={What you've built}]
One more word for telling, explaining and describing at B1.
\end{grammarlens}

\subsection*{Your turn}
\begin{itemize}
\raggedright
  \item \emph{Say it:} \emph{el alivio}
  \item \emph{Say it:} \emph{el alivio}, once more
  \item \emph{Say it:} a sentence of your own with \emph{el alivio}
  \item \emph{Recall:} say the Spanish for discouraged, then the Spanish for to cheer up, then say \emph{el alivio} again
\end{itemize}

\begin{tcolorbox}[breakable,colback=teal!4,colframe=teal!35!black,title={Before you move on}]
What does el alivio mean? (\textquotedblleft{}Relief\textquotedblright{}.) Say it once more.
\end{tcolorbox}
